\documentclass[10pt,a4paper]{amsart}
\usepackage{amsmath,amssymb,amsthm,graphicx}
\usepackage{verbatim}
\newenvironment{reference}{\comment}{\endcomment}
%\newenvironment{reference}{}{}
\newenvironment{slogan}{\comment}{\endcomment}
\newenvironment{history}{\comment}{\endcomment}

% For commutative diagrams we use Xy-pic
\usepackage[all]{xy}

% We use 2cell for 2-commutative diagrams.
\xyoption{2cell}
% Theorem environments.
%
\theoremstyle{plain}
\newtheorem{theorem}[subsection]{Theorem}
\newtheorem{proposition}[subsection]{Proposition}
\newtheorem{lemma}[subsection]{Lemma}

\theoremstyle{definition}
\newtheorem{definition}[subsection]{Definition}
\newtheorem{example}[subsection]{Example}
\newtheorem{exercise}[subsection]{Exercise}
\newtheorem{situation}[subsection]{Situation}

\theoremstyle{remark}
\newtheorem{remark}[subsection]{Remark}
\newtheorem{remarks}[subsection]{Remarks}

\numberwithin{equation}{subsection}

% Macros
%
\def\lim{\mathop{\mathrm{lim}}\nolimits}
\def\colim{\mathop{\mathrm{colim}}\nolimits}
\def\Spec{\mathop{\mathrm{Spec}}}
\def\Hom{\mathop{\mathrm{Hom}}\nolimits}
\def\Ext{\mathop{\mathrm{Ext}}\nolimits}
\def\SheafHom{\mathop{\mathcal{H}\!\mathit{om}}\nolimits}
\def\SheafExt{\mathop{\mathcal{E}\!\mathit{xt}}\nolimits}
\def\Sch{\mathit{Sch}}
\def\Mor{\mathop{\mathrm{Mor}}\nolimits}
\def\Ob{\mathop{\mathrm{Ob}}\nolimits}
\def\Sh{\mathop{\mathit{Sh}}\nolimits}
\def\NL{\mathop{N\!L}\nolimits}
\def\CH{\mathop{\mathrm{CH}}\nolimits}
\def\proetale{{pro\text{-}\acute{e}tale}}
\def\etale{{\acute{e}tale}}
\def\QCoh{\mathit{QCoh}}
\def\Ker{\mathop{\mathrm{Ker}}}
\def\Im{\mathop{\mathrm{Im}}}
\def\Coker{\mathop{\mathrm{Coker}}}
\def\Coim{\mathop{\mathrm{Coim}}}

% Boxtimes
%
\DeclareMathSymbol{\boxtimes}{\mathbin}{AMSa}{"02}

%
% Macros for moduli stacks/spaces
%
\def\QCohstack{\mathcal{QC}\!\mathit{oh}}
\def\Cohstack{\mathcal{C}\!\mathit{oh}}
\def\Spacesstack{\mathcal{S}\!\mathit{paces}}
\def\Quotfunctor{\mathrm{Quot}}
\def\Hilbfunctor{\mathrm{Hilb}}
\def\Curvesstack{\mathcal{C}\!\mathit{urves}}
\def\Polarizedstack{\mathcal{P}\!\mathit{olarized}}
\def\Complexesstack{\mathcal{C}\!\mathit{omplexes}}
% \Pic is the operator that assigns to X its picard group, usage \Pic(X)
% \Picardstack_{X/B} denotes the Picard stack of X over B
% \Picardfunctor_{X/B} denotes the Picard functor of X over B
\def\Pic{\mathop{\mathrm{Pic}}\nolimits}
\def\Picardstack{\mathcal{P}\!\mathit{ic}}
\def\Picardfunctor{\mathrm{Pic}}
\def\Deformationcategory{\mathcal{D}\!\mathit{ef}}
\title{Formally etale ring maps need not be flat}
\author{The Stacks Project Authors}
\date{stacks-project:examples.tex:local-snapshot-20261012; GNU FDL 1.2 or later}
\begin{document}
\maketitle
\noindent\textit{Editorial note.} All delivered mathematical text is the author's own, copied from the pinned Stacks Project chapter file examples.tex; no statement or proof step was rewritten, reordered or removed.\par\medskip
\noindent\textit{Editorial note.} The counted proof is the author's complete argument: Counterexample to EGA 0, Example 19.10.3(i) final sentence; companion background for the formally etale case of the structure theorem.\par\medskip
\noindent\textit{Editorial note.} The label lemma-formally-etale-not-presented names the counted statement, and the proof environment immediately following it in the source is delivered verbatim as its proof.\par\medskip
\noindent\textit{Editorial note.} The counted statement is carried by the author's own lemma environment in the chapter source.\par\medskip
\noindent\textit{Editorial note.} Revision r002: regenerated from the same spec after the r001 assembly failed at compile time (the author's citations of the letter-key EGA had no bibliography backend in r001; the stale r001 job directory is retained as evidence).\par\medskip
\noindent\textit{Editorial note.} Bibliography support: the counted context cites the letter-key EGA (the Stacks Project's global bibliography entry for Grothendieck and Dieudonne, Elements de geometrie algebrique, Springer, 1971, whose Part 0 Example 19.10.3(i) and Part IV Theorem 18.4.6(i) are the statements the example refutes and locates the error in). A matching ega.bib asset is attached so the author's own citations resolve; no statement or proof text was altered.\par\medskip
\noindent\textit{Allowed standard prerequisites (editorial).}\begin{itemize}
\item Formally etale morphism of rings: the unique lifting property against square-zero ideal extensions, namely for every ring R, every ideal I of R with I squared zero, and every commutative diagram from A to A/J and A to R into R/I, there is a unique ring map A/J to R/I making both triangles commute
\item A valuation ring of an algebraically closed non-archimedean field with J its maximal ideal provides a local ring A with a nonzero proper ideal J satisfying J squared equals J; such a J cannot be finitely generated by Nakayama
\item EGA 0, Example 19.10.3(i) claims formally etale maps out of quotients behave in a way this example refutes (this was not caught in the later errata lists of EGA)
\end{itemize}\medskip
\label{section-formally-etale-nonflat}

\noindent
In this section we give a counterexample to the final sentence in
\cite[0, Example 19.10.3(i)]{EGA} (this was not one of the items caught
in their later errata lists). Consider $A \to A/J$ for a local ring $A$
and a nonzero proper ideal $J$ such that $J^2 = J$ (so $J$ isn't finitely
generated); the valuation ring of an algebraically closed non-archimedean
field with $J$ its maximal ideal is a source of such $(A, J)$. These
non-flat quotient maps are formally \'etale. Namely, suppose given a
commutative diagram
$$
\xymatrix{
A/J \ar[r] & R/I \\
A \ar[u] \ar[r]^\varphi & R \ar[u]
}
$$
where $I$ is an ideal of the ring $R$ with $I^2 = 0$. Then $A \to R$
factors uniquely through $A/J$ because
$$
\varphi(J) = \varphi(J^2) \subset (\varphi(J)A)^2 \subset I^2 = 0.
$$
Hence this also provides a counterexample to the formally \'etale case
of the ``structure theorem'' for locally finite type and formally \'etale
morphisms in \cite[IV, Theorem 18.4.6(i)]{EGA}
(but not a counterexample to part (ii), which
is what people actually use in practice). The error in the proof of
the latter is that the very last step of the proof is to invoke the incorrect
\cite[0, Example 19.3.10(i)]{EGA},
which is how the counterexample just mentioned creeps in.


\begin{lemma}
\label{lemma-formally-etale-not-presented}
There exist formally \'etale nonflat ring maps.
\end{lemma}

\begin{proof}
See discussion above.
\end{proof}

\bibliographystyle{plain}
\bibliography{ega}
\end{document}
